\section{The cost of testing many lists}\label{sec:list}

Two distinctness tests suffice to check that records have reached their
specified destinations. This gives a common starting point for the two
cases of Theorem~\ref{thm:many}. In a field without ports, the possible
destinations force many wires to have large total length in a suitable
spectral space. In a finite field with ports, cheap distinctness tests
would make it possible to compress a family of tensors with a known
lower bound.

\subsection{Checking a tagged permutation}

Write \(h=w-1\), \(N=2^h\), and index records by the \(h\)-bit words.
The tagged-permutation tensor \(\operatorname{Tag}(h,h)\) of
\cite[Definition~\ext{G}{def:tag}]{oaG} takes
\(N\) records \((j_i,x_i)\) and an ordered list of \(N\) payloads
\(y_j\), with \(h\) bits per tag or payload. Its value is one when
\(i\mapsto j_i\) is a permutation and \(x_i=y_{j_i}\) for every \(i\).
Its other entries are zero.

The first \(D_w\) test checks the tags. Encode each \(j_i\) as
\((0,j_i,0^{h+1})\) and fill the remaining \(N\) positions with the
distinct words \((1,j,0^{h+1})\). All \(2N\) words have length \(2w\),
and the test passes precisely when the tags form a permutation \(\pi\).

The second \(D_w\) test checks the payloads. Its first \(N\) words are
\((j_i,0,x_i,0)\). At the position for \(j\) in the second half, sum
over two choices, the reserve word \((j,1,0^h,0)\) with weight \(1\)
and the comparison word \((j,0,y_j,0)\) with weight \(-1\).
For tags forming \(\pi\), the only possible collision at address \(j\)
is between the comparison word and the word from \(i=\pi^{-1}(j)\).
Thus the signed sum factors as
\[
     \prod_j \bigl(1-\mathbf 1_{\{x_{\pi^{-1}(j)}\ne y_j\}}\bigr)
       =\prod_j\mathbf 1_{\{x_{\pi^{-1}(j)}=y_j\}} .
\]
This identity holds in every characteristic. A summed selector bit for
each address controls its word and its weight. Copying the tags,
selecting the words, and pinning the constants cost \(O(Nh)\) trivalent
volume. Consequently, if \(D_w^{\otimes t}\) has a realization of
volume \(V\), two uses of that realization give a network for
\(\operatorname{Tag}(h,h)^{\otimes t}\) of volume at most
\(2V+O(tNh)\), including when the realization shares work among lists.

\subsection{Fields without ports}

A class of Hadamard spaces has \emph{ports} if it admits regular
graphs of unbounded order and common bounded degree with the following
property. For every placement of their vertices into any space of the
class, the mean squared distance of two independent uniform vertices
is at most a fixed constant times the mean squared edge length.
A field has ports when its Hadamard spectral spaces have this property
\cite[Definition~\ext{G}{def:ports}]{oaG}.

\begin{lemma}[Uniform destinations]\label{lem:uniform-demand}
Let the field have no ports. Consider a fixed trivalent topology with
\(r\) input bits and \(r\) output bits. Suppose that for every permutation
in the support of a distribution, some choice of vertex tables realizes
the \(r\) independent wires joining each input to its assigned output.
If each input has a uniform destination on a set of \(q\ge2\) outputs,
the topology has size \(\Omega(r\sqrt{\log q})\). The destination sets
may depend on the input.
\end{lemma}

\begin{proof}
Give each exterior bit its own equality leaf. Connect the resulting
components in a path by adding edges to the graph used to measure
distances. This increases the maximum degree by at most two and gives
\(m\) edges, with \(m\) at most a constant times the network size.
Let \(\alpha\) be the law of a uniform oriented edge, and let \(\mu\)
be the joint law of a uniform input and its random destination.
For any placement in a spectral space with distance \(d\),
\[
                    r\,\mathbb E_\mu d\ \le\ m\,\mathbb E_\alpha d .
\]
Indeed, for each permutation the network tensors locally contract
the product of its internal bit wires to the prescribed exterior
wires. Monotonicity and additivity of prices bound the sum of the
destination distances by the sum of the internal edge lengths.
Average over the permutations and include the extra edges on the right.

We need a placement where many destinations are far away at once.
The class of Hadamard spectral spaces is closed under isometry,
rescaling, closed convex subspaces, and pointed ultralimits
(Section~\ref{sec:prices-background}). The product theorem
\cite[Theorem~\ext{P}{prod:main}]{oaP} gives a subclass with the same
properties and no ports, closed also under finite \(\ell_2\) products.
Let \(\mathcal K\) be the squared distance tables \(D=d^2\) realized
in this subclass with \(\mathbb E_\alpha D\le1\). Products of rescaled
placements give convex combinations of tables. Connectivity bounds
every entry, and pointed ultralimits realize coordinatewise limits.
Hence \(\mathcal K\) is convex and compact.

For a probability law \(\nu\le 2\mu\) on ordered pairs, let
\(\beta_\nu=\max_{D\in\mathcal K}\mathbb E_\nu D\). By rescaling,
\(\mathbb E_\nu d^2\le\beta_\nu\mathbb E_\alpha d^2\) for every
placement in the subclass. Write \(X,Y\) for the endpoints under
\(\nu\). The graph entropy bound \cite[Corollary~\ext{P}{cor:graph-entropy}]{oaP}
and the bounded degree of our graph give
\(H_\nu(Y\mid X)\le C(1+\beta_\nu)\), with entropy in bits.
Every pair has \(\nu\)-mass at most \(2/(rq)\), so
\[
    H_\nu(Y\mid X)=H_\nu(X,Y)-H_\nu(X)
               \ge \log_2(rq/2)-\log_2 r=\log_2(q/2).
\]
It follows that \(\beta_\nu\ge c\log q\) for all sufficiently large
\(q\), uniformly in \(\nu\).

Bilinear minimax on \(\mathcal K\) and the compact convex set of
laws \(\nu\le2\mu\) gives a single \(D=d^2\in\mathcal K\) for which
\(\mathbb E_\nu D\ge c\log q\) for every such law. If the pairs with
\(D<c\log q\) had \(\mu\)-mass at least \(1/2\), conditioning \(\mu\)
on those pairs would violate this bound. Hence
\(\mathbb E_\mu d\ge \tfrac12\sqrt{c\log q}\), while
\(\mathbb E_\alpha d\le1\) by Cauchy--Schwarz. The price inequality
gives the result. For bounded \(q\), counting exterior bits suffices.
\end{proof}

In a network for \(\operatorname{Tag}(h,h)^{\otimes t}\), pin each
list's tags to a uniform permutation. The topology realizes a permutation
of \(r=tNh\) payload bits. Each input bit has \(q=N\) possible
destinations, at its own bit position in its own list, all equally
likely. By Lemma~\ref{lem:uniform-demand},
\(2V+O(tNh)\ge c\,tNh\sqrt h\). Since \(n_w=4N(h+1)\), this proves
\eqref{eq:many} for any fixed \(0<\sigma<1/2\) and all sufficiently
large \(w\), uniformly in \(t\).

\subsection{Finite fields with ports}

Fix a finite field with ports. We use the signed counting tensor
\(U_\ell^*\) on \(\ell\) exterior bits from the circuit lower-bound
paper. Its entries count the signed computation paths of a polynomially
clocked universal machine, reduced in the field. The polynomial-size
verifier gives an exact trivalent network of size \(\ell^{O(1)}\) for
one copy. Sum over its guess strings, requiring the unused bits on
each computation path to be zero, and weight the outcome by its sign. The exact batch lower bound
\cite[Theorem~\ext{P}{thm:batches}]{oaP} gives \(\delta,\xi>0\)
such that, for all sufficiently large \(\ell\),
\[
       \operatorname{vol}_k\bigl((U_\ell^*)^{\otimes j}\bigr)
              >j\ell^{1+\delta}
       \qquad (1\le j\le 2^{\ell^\xi}).
\]

The descent with tagged permutations
\cite[Proposition~\ext{G}{prop:oracle-descent}]{oaG} supplies a
constant \(a>0\). For all sufficiently large \(h\), given a trivalent
network of size \(s\) in the range \(N\le s\le N^4\), with \(e\)
exterior bits, it constructs a network for \(d\) copies of its tensor
using trivalent gates and \(\operatorname{Tag}(h,h)\) tests. Its size
\(S\), charging every incident bit including those at tests, satisfies
\[
            S\le C d\bigl(sh^{-a}+e(\log h)^C\bigr),
           \qquad \log_2 d=O\bigl((\log h)^2\bigr).
\]
These bounds are uniform in the network. The normalization required
for descent attaches equality leaves and subdivides self-loops. It
costs only a constant factor, and the cited bound holds on every fixed
polynomial size range.

Choose \(b>\max\{2,2/\xi\}\) and
\(0<\sigma<\min\{a/4,b\delta/4,1/4\}\). Suppose there are arbitrarily
large \(w=h+1\) for which some \(1\le t\le 2N\) admits a realization
of \(D_w^{\otimes t}\) of volume at most \(tn_w w^\sigma\).
The two-list construction gives a network for \(t\) tagged-permutation
tests of size \(O(tNh^{1+\sigma})\).

Group the tests in a descent network into batches of \(t\) and
substitute this realization for each batch. In the last batch pin
unused tests to the identity permutation with zero payloads, where
each has value one. Since the substituted network realizes the full
product tensor, the replacement stays exact under every contraction
of the surrounding wires, including wires joining tests in one batch.
Each full batch contributes \(3tNh\) incident bits to \(S\).
Paying separately for the last batch therefore gives trivalent size
\[
                     O\bigl(h^\sigma S+tNh^{1+\sigma}\bigr).
\]

Set \(\ell=\lceil h^b\rceil\) and begin with \(j=N^2\) copies of
the polynomial-size network for \(U_\ell^*\). At each descent let
\(x=s/(j\ell)\) be the size per exterior bit. The new batch has \(dj\)
copies. Its size per bit, for all sufficiently large \(h\), satisfies
\[
 \begin{split}
 x'&\le C'\left(h^{-a+\sigma}x+
        h^\sigma(\log h)^C+\frac{tNh^{1+\sigma}}{dj\ell}\right)\\
   &\le h^{-a/2}x+h^{2\sigma}.
 \end{split}
\]
Here \(j\ge N^2\), \(d\ge1\), and \(t\le 2N\), so the fraction is
\(O(h^{1+\sigma-b})\). This is the reason for starting with \(N^2\)
copies.

Initially \(x\le h^A\) for some fixed \(A\). Choose a fixed integer
\(L>2A/a\). At every one of these \(L\) steps, the product of the
replication factors is \(N^{o(1)}\). Inductively, \(j=N^{2+o(1)}\)
and \(x\) stays polynomial in \(h\), so the current network and its
normalization have size in the descent range. Iterating the last
inequality gives \(x=O(h^{2\sigma})\). Now \(j\le 2^{\ell^\xi}\)
because \(\log_2 j=(2+o(1))h\), while
\(h^{2\sigma}=o(\ell^\delta)\) by the choice of \(\sigma\).
The final trivalent volume is therefore \(o(j\ell^{1+\delta})\),
contradicting the exact batch lower bound. This proves
Theorem~\ref{thm:many} in the remaining case.
